Cross-validation (CV) can underestimate spatial prediction error when validation observations are closer to training data than the intended prediction locations. We examine whether accounting for spatial correlation improves the choice of CV scheme beyond matching distances alone. We define $J_2$ as the difference between validation and target locations in their average squared correlation with the nearest training observation. For Gaussian prediction using one neighbour, with known mean and covariance, multiplying $J_2$ by response variance gives target mean squared error (MSE) minus expected CV MSE. Using additional training observations can further reduce prediction error, but this reduction can differ between validation and target locations. We derive this difference and show that it can produce CV bias even when the nearest-distance distributions match exactly. Among four CV schemes, selection by $|J_2|$ reduced mean absolute expected bias by \BiasReduction\% relative to selection by the Wasserstein distance $W$ across 5,120 configurations with known covariance. In 1,920 datasets with estimated covariance, the reduction in mean absolute error (MAE) of the observed CV estimate was \NestedReduction\%. Almost all of this gain came from one setting combining uniform sampling, rapidly decaying spatial correlation, low observation noise and a prediction area extending beyond the sampled region. A ten-fold check gave a 1.7\% reduction, but the Monte Carlo interval for the MAE difference included zero. Further simulations show that information supplied by predictors can change CV bias while response covariance remains fixed. In a soil-cadmium application, both diagnostics select kNNDM, but random CV is closer to the observed prediction error at 100 reserved sites. Nearest-neighbour correlation alone cannot generally determine CV bias when prediction uses several neighbours or predictors.
